\setcounter{footnote}{90}

\printedpage{98}
% head: A. Grothendieck : Dérivateurs, Ch. XII, Édition des Universités de Montpellier II et Paris VII
% foot: Version 7 décembre 2002
\[
\begin{tikzcd}[row sep=small, column sep=small]
 & X_1^\circ \arrow[dr,"i^{\prime}_{1}"] & & \\
\widehat{Q^\circ} = X_0^\circ \arrow[ur,"\alpha_1^\circ"] \arrow[dr,"\alpha_2^\circ"'] \arrow[rr,dotted] & & \mathcal{X}_{Q^\circ} \arrow[r] & Y^\circ, \\
 & X_2^\circ \arrow[ur,"i^{\prime}_{2}"'] & &
\end{tikzcd}
\]
\[
\begin{tikzcd}[row sep=small, column sep=small]
 & X_1 \arrow[dr] & & \\
\widehat{Q}\tld = X_0 \arrow[ur] \arrow[dr] \arrow[rr,dotted] & & \tilde{\mathcal{X}}_{Q} \arrow[r] & Y, \\
 & X_2 \arrow[ur] & &
\end{tikzcd}
\]
\[
\begin{tikzcd}[row sep=small, column sep=small]
 & X_1^\circ \arrow[dr] & & \\
\widehat{Q^\circ}\tld = X_0^\circ \arrow[ur] \arrow[dr] \arrow[rr,dotted] & & \tilde{\mathcal{X}}_{Q^\circ} \arrow[r] & Y^\circ, \\
 & X_2^\circ \arrow[ur] & &
\end{tikzcd}
\]
(\footnote{\textbf{NB} We should still complete these (noncommutative) diagrams with the dotted arrows \dots}) \textit{are isomorphic in} $\Hot_W(\Gamma)$\textit{. More precisely still, we have the isomorphisms in} $\uHom(\Gamma^\circ,\Cat)$
\[
(\widehat{Q}\tld)^\circ \simeq \widehat{Q^\circ},
\]
\textit{whence, applying it to $Q^\circ$ instead of $Q$,}
\[
(\widehat{Q^\circ}\tld)^\circ \simeq \widehat{Q},
\]
\textit{which implies (via the canonical isomorphisms valid in} $\HOT_W(S)$ \textit{for every object $S$ in} $\Cat$ \textit{and every} $X_* \in \mathrm{Ob}\,\uHom(S,\Cat)$\textit{,} $X_* \simeq X_*^\circ$ (\footnote{via the functor $\Sigma$ used above \dots})\textit{), that we have the canonical isomorphisms}
\[
\left.\begin{aligned}
\widehat{Q} &\simeq \widehat{Q^\circ}\tld \\
\widehat{Q^\circ} &\simeq \widehat{Q}\tld
\end{aligned}\right\}
\qquad \text{\textit{in} } \HOT_W(\Gamma),
\]
\textit{but also directly}
\[
\widehat{Q} \simeq \widehat{Q^\circ}, \qquad \widehat{Q}\tld \simeq \widehat{Q^\circ}\tld \qquad \text{\textit{in} } \HOT_W(\Gamma)
\]
\msp{page 143}\textit{(and it remains to check that we thus have a} transitive \textit{system of canonical isomorphisms in} $\HOT_W(\Gamma)$\textit{, between the four objects considered, which must be essentially sorital \dots).}

\subsection*{Remarks.}

\begin{enumerate}
\item[1)] What we have just said in the case of diagrams in Cat of type $\Phi^\circ = \Psi$, seems still to make sense for diagrams of arbitrary type
\[
X_* : S^\circ \longrightarrow \Cat, \qquad s \longmapsto X_s
\]
\printedpage{99}
% head: A. Grothendieck : Dérivateurs, Ch. XII, Édition des Universités de Montpellier II et Paris VII
% foot: Version 7 décembre 2002
($S$ an object of Cat). We can integrate it in a smooth way, $\int X_* = \int_{\mathrm{lis}} X_*$, obtaining a fibered category (split) over $S$, or in a proper way, $\int_{\mathrm{pr}} X_*$, obtaining a split cofibered category over $S^\circ$. We can also consider $X_*^\circ$, the composite of $X_*$ with the automorphism $Y \longmapsto Y^\circ$ in Cat, and integrate it in either a proper or a smooth way. In all these cases, we find a diagram of type $\widehat{S}^\circ$, where $\widehat{S}$ is obtained from $S$ by adding an element which maps by a ``canonical'' (given) arrow to each of the others (\textit{without} commutation relations!) A
\msp{page 144}\textit{direct system} of morphisms in Cat
\[
f \qquad X_s \xrightarrow{f_s} Y
\]
then gives a morphism
\[
\int_{\mathrm{lis}} X_* \xrightarrow{\varphi_Q} Y,
\]
where $Q$ denotes the diagram, of type $\widehat{S}^\circ$, made up of the $X_s$, $Y$, the $u^* : X_t \longrightarrow X_s$ (for $u : s \longrightarrow t$ in $S$) and the $f_s$. Likewise we will have
\[
\int_{\mathrm{prop}} X_* \xrightarrow{\tilde{\varphi}_Q} Y,
\]
and if
\[
\Gamma = (\widehat{S} + \text{initial object}),
\]
we find diagrams $\widehat{Q}$, $\widehat{Q}\tld$ of type $\Gamma^\circ$ in Cat. Unless I am mistaken, the relations between $\widehat{Q}$, $\widehat{Q}\tld$, $\widehat{Q^\circ}$, $\widehat{Q^\circ}\tld$ are the same as those made explicit in proposition 14 above. I am not going to dwell on this with the requisite care, unless needed. We will say that a diagram $Q$ defines $Y$ as a $\varinjlim\limits_{I}$ that is $W$-homotopic \texttt{[where} $I = S^\circ$\texttt{]} if $\varphi_Q \in W$. It amounts to the same thing to say that $\varphi_{Q^\circ} \in W$, or that $\tilde{\varphi}_Q \in W$ or that $\tilde{\varphi}_{Q^\circ} \in W$ \dots
\msp{page 145}

\item[2)] The ``universal'' case, for
\begin{tikzcd}[baseline=(\tikzcdmatrixname-2-1.base), row sep=small, column sep=small]
 & X_1 \\
X_0 \arrow[ur] \arrow[dr] & \\
 & X_2
\end{tikzcd}
given, of a commutative square $Q = {}$%
\begin{tikzcd}[baseline=(\tikzcdmatrixname-2-1.base), row sep=tiny, column sep=small]
 & X_1 \arrow[dr] & \\
X_0 \arrow[ur] \arrow[dr] & & Y \\
 & X_2 \arrow[ur] &
\end{tikzcd}
in Cat, giving rise to
\[
\begin{tikzcd}[row sep=small, column sep=small]
 & & & X_1 & & & \\
\varphi_Q : &[1em] \left(\vrule height 6ex depth 6ex width 0pt\right. &[-1.5em] \int X_0 \arrow[ur] \arrow[dr] & & \left.\vrule height 6ex depth 6ex width 0pt\right) &[2em] \longrightarrow &[2em] Y, \\
 & & & X_2 & & &
\end{tikzcd}
\]
is that of $Y = X_1 \amalg_{X_0} X_2$, giving rise to % [sic] naissace -> naissance, letter n omitted
\printedpage{100}
% head: A. Grothendieck : Dérivateurs, Ch. XII, Édition des Universités de Montpellier II et Paris VII
% foot: Version 7 décembre 2002
\[
\begin{tikzcd}[row sep=small, column sep=small]
 & X_1 \arrow[dr] & \\
Q_0 = Q(X_*) = X_0 \arrow[ur] \arrow[dr] & & X_1 \amalg_{X_0} X_2. \\
 & X_2 \arrow[ur] &
\end{tikzcd}
\]

The general case reduces to this one, for we will have a commutative diagram
\[
\begin{tikzcd}[row sep=small, column sep=small]
 & & X_1 & & & \\
\left(\vrule height 6ex depth 6ex width 0pt\right. &[-1.5em] \int X_0 \arrow[ur] \arrow[dr] & & \left.\vrule height 6ex depth 6ex width 0pt\right) \arrow[rr,"{\varphi_{Q_0}=\varphi_{Q(X_*)}}"] \arrow[ddrr,"\varphi_Q"] &[4em] \hphantom{MMMM} &[4em] X_1 \amalg_{X_0} X_2 \arrow[dd,"\alpha_Q"] \\
 & & X_2 & & & \\
 & & & & & Y.
\end{tikzcd}
\]
\end{enumerate}

But there must be interesting cases where $\varphi_{Q_0}$ is \textit{not} in $W$, i.e. the Cat-cocartesian square
\begin{tikzcd}[baseline=(\tikzcdmatrixname-2-1.base), row sep=tiny, column sep=small]
 & X_1 \arrow[dr] & \\
X_0 \arrow[ur] \arrow[dr] & & X_1 \amalg_{X_0} X_2 \\
 & X_2 \arrow[ur] &
\end{tikzcd}
is not $W$-homotopically cocartesian, and where nevertheless we can find a commutative square $Q$
\msp{page 146}completing $X_*$, in such a way that this square is $W$-homotopically commutative \texttt{[rather cocartesian]}, i.e. that $\varphi_Q = \alpha_Q \varphi_{Q_0}$ is in $W$. This is all the more plausible since in Cat, amalgamated sums tend not to have good homotopy properties, even if one of the arrows $\alpha_1$, $\alpha_2$ is \texttt{[a]} monomorphism (or even an open immersion), contrary to what happens in $\boldsymbol{\Delta}^\wedge$, and more generally in categories of the form $Z^\wedge$ ($Z$ in Cat).

\subsection*{3.4\quad Criteria for homotopically cocartesian squares.}

The criteria I have obtained are all contained in the following statement:

\textbf{Theorem 4} (\footnote{\textbf{NB} In this theorem, and corollaries 1, 2, 3, we need only the \textit{weak form} of W(7 bis), namely $i_0 \in W \Longrightarrow i \in W$ (cf. p. 79) (Loc(6$'$a) in the notation of XVI).},\footnote{\textbf{NB} All in all, the theorem (forward part) holds without any hypothesis W(7 bis) or W(7 bis faible) on $W$, and the same is true of its corollaries 1 to 3 (to the exclusion of corollaries 4, 5, which do use W(7 bis) in its full strength).}). \textit{Let us consider a commutative square diagram in} Cat
\[
Q \quad = \quad
\begin{tikzcd}
X_0 \arrow[r,"\alpha_1"] \arrow[d,"\alpha_2"'] \arrow[dr,dotted,"f_0"] & X_1 \arrow[d,"f_1"] \\
X_2 \arrow[r,"f_2"'] & Y,
\end{tikzcd}
\tag*{$(*)$}
\]
\printedpage{101}
% head: A. Grothendieck : Dérivateurs, Ch. XII, Édition des Universités de Montpellier II et Paris VII
% foot: Version 7 décembre 2002
\textit{and consider the canonical morphism}
\[
\begin{tikzcd}[row sep=small, column sep=small]
 & & X_1 & & & \\
\varphi_Q : \quad \mathcal{X} = \Bigl(\int &[-1.2em] X_0 \arrow[ur,"\alpha_1"] \arrow[dr,"\alpha_2"'] & & \Bigr) &[2em] \longrightarrow &[2em] Y, \\
 & & X_2 & & &
\end{tikzcd}
\]
\msp{page 147}\textit{the integral being the integral which is} smooth. \textit{(There is a dual statement concerning the proper integral of}\newline
\begin{tikzcd}[baseline=(\tikzcdmatrixname-2-1.base), row sep=small, column sep=small]
 & X_1 \\
X_0 \arrow[ur] \arrow[dr] & \\
 & X_2
\end{tikzcd}
\textit{.)}

\begin{enumerate}
\item[\textit{a)}] \textit{For $\varphi_Q$ to be $W$-smooth, it is necessary and sufficient that $f_1$, $f_2$ and $f_0 = f_1\alpha_1 = f_2\alpha_2$ be so.} (Assuming W(7 bis) (\footnote{For the forward part ($f_0$, $f_1$, $f_2$ smooth $\Longrightarrow \varphi_Q$ smooth), we need only W(7 bis) (\footnotemark), weak form. It is only for the converse (which has never yet been of use to me) that we need W(7 bis), in its full strength.}\footnotetext{We do not even need W(7 bis) in weak form, cf. annotation on page 156. We use a restrictive hypothesis on $W$ only in the converse ($\varphi_Q$ smooth $\Longrightarrow f_0$, $f_1$, $f_2$ smooth).})).

\item[\textit{b)}] \textit{The fiber of $\varphi_Q$ at an element $y$ of $Y$ is isomorphic to}
\[
\begin{tikzcd}[row sep=small, column sep=small]
 & & X_{1\,y} & \\
\int &[-1.2em] X_{0\,y} \arrow[ur] \arrow[dr] & &[2em] (\simeq \mathcal{X}_y). \\
 & & X_{2\,y} &
\end{tikzcd}
\tag*{$(**)$}
\]
\end{enumerate}

\textbf{Corollary 1.} (\footnote{Valid without any hypothesis on $W$ (other than Loc(1,2,3 bis)).}) \textit{Suppose $f_1$, $f_2$ and $f_1\alpha_1 = f_2\alpha_2$ $W$-smooth, and the integrated categories $(**)$ $W$-aspheric for every $y \in \mathrm{Ob}\,Y$. (NB This is still the smooth integral, but it does not really matter, for we have seen (cf. the upper part of the diagram p. 140, which implies that $\mathcal{X}$} \texttt{[is]} \textit{$W$-aspheric if and only if $\tilde{X}$ is) that one is $W$-aspheric if and only if the other is.) Then $\varphi_Q \in W$, i.e. the square $(*)$ is $W$-homotopically cartesian} \texttt{[rather cocartesian]}\textit{. By duality, same conclusion if we suppose $f_1$, $f_2$, $f_1\alpha_1 = f_2\alpha_2$} proper \textit{instead of smooth. (For $Q$ $W$-homotopically trivial} \texttt{[rather cocartesian]}\textit{, if and only if $Q^\circ$ is, by} \texttt{[the]} \textit{prop. 14.)}

\msp{page 148}Indeed, a morphism in Cat which is $W$-smooth and has $W$-aspheric fibers is in $W$ (and it is even a universal $W$-equivalence \dots).

\textbf{Corollary 2}. \textit{We assume that $W$ satisfies} W(7 bis) (cocartesian square axiom) (\footnote{unnecessary.}).
\printedpage{102}
% head: A. Grothendieck : Dérivateurs, Ch. XII, Édition des Universités de Montpellier II et Paris VII
% foot: Version 7 décembre 2002
\textit{Suppose the square $Q$ cocartesian in} Cat (\footnote{NB We do not really use that $Q$ is cocartesian, rather that $Q$ is \textit{cartesian,} and $f_2$ (hence also $\alpha_1$) an open immersion (resp. closed), and that for $y \in Y\backslash X_2$, $X_{1\,y}$ is $W$-aspheric. (In the case where $Q$ is cocartesian, $X_{1\,y}$ is $\simeq e$.)})\textit{, and that $\alpha_1$ is an open immersion (resp. closed) and $f_1$ smooth (resp. proper). Then $Q$ is $W$-homotopically cocartesian.} (\footnote{Here we need W(7 bis), forward form, i.e. Loc(6$'$a) (notation of XVI).})

It suffices to establish the non-resp'ed result, the other will follow from it by duality (since $Q$ is $W$-cocartesian if and only if $Q^\circ$ is). By hypothesis $f_1$ is $W$-smooth, on the other hand $\alpha_1$ and $f_2$ are open immersions, hence $W$-smooth, hence $f_1\alpha_1$ is $W$-smooth too. (Besides, $\alpha_2$ too, for the square is also cartesian, hence $\alpha_2$ obtained from $f_1$ (which is $W$-smooth by hypothesis) by base change $X_2 \longrightarrow Y$.) Thus $\varphi_Q$ is \texttt{[}$W$\texttt{-]}smooth, and it suffices to see that its fibers are $W$-aspheric. The calculation $(**)$ gives, for $y$ in $Y\backslash X_2$,
\[
\begin{tikzcd}[row sep=small, column sep=small]
 & & e & \\
\mathcal{X}_y \simeq \quad \int &[-1.2em] \emptyset \arrow[ur] \arrow[dr] & &[1em] \simeq e, \qquad \text{O.K.} \\
 & & \emptyset &
\end{tikzcd}
\]
And for $y$ in $X_2$
\[
\begin{tikzcd}[row sep=small, column sep=small]
 & & X_{1\,y} & \\
\mathcal{X}_y \simeq \int &[-1.2em] X_{0\,y} \arrow[ur,"\sim"] \arrow[dr] & &[2em] . \\
 & & e &
\end{tikzcd}
\]
\msp{page 149}Since $X_{0\,y} \longrightarrow X_{1\,y}$ is an isomorphism, hence is in $W$, it follows that $e \longrightarrow \mathcal{X}_y$ is too (prop. 13, p. 125), assuming W(7 bis), hence $\mathcal{X}_y$ is $W$-aspheric, O.K. (\footnote{We need only the forward form of W(7 bis), i.e. (notation of XVI) Loc(6$'$a).}).

\textbf{Corollary 3.} (Assuming W(7 bis) (\footnote{unnecessary.}).) \textit{Let $X$ be in} Cat\textit{, $U_1$, $U_2$ two open subcategories of $X$ such that $X = U_1 \cup U_2$, $U_0 = U_1 \cap U_2$ the open intersection. Then the square of open immersions}
\[
\begin{tikzcd}[row sep=large, column sep=large]
U_0 \arrow[r,hook] \arrow[d,hook] & U_1 \arrow[d,hook] \\
U_2 \arrow[r,hook] & X
\end{tikzcd}
\]
\textit{is $W$-cocartesian. (Dual statement for two closed subcategories $T_1$, $T_2$ such that $T_1 \cup T_2 = X$, with $T_0 = T_1 \cap T_2$.)}

We apply corollary 2. We know that the square is cocartesian in Cat (\footnote{We do not even need to know this, cf. marginal annotation to corollary 2.}), moreover $U_1 \lhook\joinrel\longrightarrow X$ is indeed $W$-smooth, being an open immersion.